\documentclass[aspectratio=169]{beamer}

\usetheme{Madrid}
\usecolortheme{dove}

\usepackage[T1]{fontenc}
\usepackage[utf8]{inputenc}
\usepackage{array}
\usepackage{booktabs}
\usepackage{graphicx}
\usepackage{hyperref}

% Muted forensic palette: dark slate, steel blue, soft amber.
\definecolor{sbslate}{RGB}{35,45,56}
\definecolor{sbblue}{RGB}{67,97,125}
\definecolor{sblightblue}{RGB}{94,121,145}
\definecolor{sbamber}{RGB}{156,124,72}
\definecolor{sbpaper}{RGB}{246,247,248}

\setbeamercolor{title}{fg=sbslate}
\setbeamercolor{subtitle}{fg=sbblue}
\setbeamercolor{frametitle}{bg=sbslate,fg=white}
\setbeamercolor{structure}{fg=sbblue}
\setbeamercolor{block title}{bg=sbblue,fg=white}
\setbeamercolor{block body}{bg=sbpaper,fg=black}
\setbeamercolor{alerted text}{fg=sbamber}
\setbeamerfont{block title}{series=\bfseries}
\setbeamertemplate{navigation symbols}{}
\setbeamertemplate{footline}[frame number]
\setbeamersize{text margin left=0.75cm,text margin right=0.75cm}

\title[Shattered Bytes]{Shattered Bytes}
\subtitle{A Manual Data Carving Serious Game for Computer Forensics Education}
\author{Simone Romano}
\institute{Computer Forensics and Cyber Crime Analysis}
\begin{document}

% 1
\begin{frame}
  \titlepage
\end{frame}

% 2
\begin{frame}{Design Focus}
  \begin{block}{What the game trains}
    \emph{Shattered Bytes} focuses on the complete \textbf{manual data carving workflow}:
    detecting meaningful byte patterns, rejecting false leads, reconstructing fragmented
    evidence, and reporting only what can be defended from the recovered data.
  \end{block}

  \vspace{0.2cm}
  \begin{itemize}
    \item The player must combine \textbf{signature analysis}, \textbf{offset reasoning},
      and \textbf{fragment reconstruction}.
    \item The campaign avoids a pure ``find the header'' exercise by adding decoys,
      partial recovery, XOR obfuscation, and partition-based navigation.
    \item Recent gameplay additions include a drag-and-drop \textbf{evidence intake}
      sequence and a final \textbf{incident timeline board}.
    \item The final output is both the recovered evidence and a transparent
      \textbf{competency report} on the player's investigative method.
  \end{itemize}
\end{frame}

% 3
\begin{frame}{Motivation}
  \begin{columns}[T,onlytextwidth]
    \begin{column}{0.49\textwidth}
      \begin{block}{Teaching problem}
        Data carving can look like simple \textbf{signature matching}, but real recovery
        also requires boundary validation, fragment ordering, decoy rejection, and careful
        reporting.
      \end{block}
    \end{column}
    \begin{column}{0.49\textwidth}
      \begin{block}{Design answer}
        The player works directly on \textbf{raw bytes}, selects exact ranges,
        reconstructs evidence, verifies the result, and submits a forensic conclusion.
      \end{block}
    \end{column}
  \end{columns}

  \vspace{0.3cm}
  The game turns a low-level forensic task into a sequence of \textbf{explicit investigative decisions}.
\end{frame}

% 4
\begin{frame}{Narrative Frame}
  \begin{block}{Operation Shattered Meridian}
    A regional hospital network has been hit by the Night Meridian ransomware cell.
    The live breach is contained, but the payment channel, mule identities, ATM cash-out,
    fraudulent transfer, ransomware payload, and exfiltration path are still hidden in
    seized devices.
  \end{block}

  \vspace{0.2cm}
  \begin{itemize}
    \item \textbf{Player role}: junior forensic analyst assigned to Agent Root.
    \item \textbf{Mission}: reconstruct six digital evidence items from raw dumps.
    \item \textbf{Final result}: one defensible incident chain.
  \end{itemize}
\end{frame}

% 5
\begin{frame}{Core Gameplay Loop}
  \begin{enumerate}
    \item \textbf{Prepare} the evidence workflow: sealed drive, write blocker,
      forensic workstation, hash verification.
    \item \textbf{Identify} the forensic target from the mission briefing.
    \item \textbf{Inspect} the dump with \texttt{search}, \texttt{go}, \texttt{select}, and \texttt{info}.
    \item \textbf{Stash} exact byte ranges in the workbench.
    \item \textbf{Compose} fragments, apply XOR when needed, and carve the artifact.
    \item \textbf{Verify} the recovered image or text payload in the asset viewer.
    \item \textbf{Report}: \texttt{recovered}, \texttt{partial}, or \texttt{inconclusive}.
  \end{enumerate}

  \vspace{0.2cm}
  \begin{block}{Design principle}
    The player is rewarded for a disciplined forensic workflow: observe, validate,
    reconstruct, and \textbf{report only what the bytes support}.
  \end{block}
\end{frame}

% 6
\begin{frame}{Campaign Progression}
  \small
  \setlength{\tabcolsep}{3pt}
  \begin{tabular}{@{}>{\bfseries}p{0.07\textwidth}p{0.25\textwidth}p{0.32\textwidth}p{0.2\textwidth}@{}}
    \toprule
    Act & Topic & Main forensic skill & Outcome \\
    \midrule
    1 & False PNG lead & Header/footer validation & Full recovery \\
    2 & Fragmented PNG & Descriptor-based reconstruction & Full recovery \\
    3 & Multi-signature triage & Relevance over recognisable files & Full recovery \\
    4 & MBR partition & Little-Endian LBA to byte offset & Full recovery \\
    5 & Obscured tail & XOR known-plaintext and partial recovery & Partial report \\
    6 & Ransomware aftermath & Multi-fragment XOR recovery & Full recovery \\
    \bottomrule
  \end{tabular}
\end{frame}

% 7
\begin{frame}{Interface and Tools}
  \begin{itemize}
    \item \textbf{Hex Editor}: byte-level inspection and range selection.
    \item \textbf{Terminal}: forensic commands such as \texttt{search}, \texttt{go},
      \texttt{select}, \texttt{entropy}, \texttt{xorcalc}, \texttt{xor}, and \texttt{report}.
    \item \textbf{Workbench}: fragment storage, ordering, XOR transformation, carving.
    \item \textbf{Asset Viewer}: visual/text verification of recovered evidence.
    \item \textbf{Evidence Intake}: drag-and-drop acquisition procedure before Act 1.
    \item \textbf{Timeline Board}: final drag-and-drop correlation of all recovered artefacts.
    \item \textbf{Objectives Panel}: explicit investigation goals and manual hints.
  \end{itemize}
\end{frame}

% 8
\begin{frame}{Assessment Model}
  \begin{block}{Score formula}
    \small
    \texttt{score = maxScore - 15*hints - 25*badSelections - 30*extraCarves
    - 40*wrongReports - 15*procedureMistakes + timeBonus}
  \end{block}

  \vspace{0.15cm}
  \begin{itemize}
    \item \textbf{Bad selections}: decoys, unsupported ranges, overlaps, or extra bytes.
    \item \textbf{Extra carves}: carving before validating boundaries and order.
    \item \textbf{Wrong reports}: overclaiming partial evidence or misclassifying recovery.
    \item \textbf{Procedure mistakes}: failed post-action forensic decisions.
    \item \textbf{Final mastery}: signature recognition, offset precision, fragment
      reconstruction, trial-and-error discipline, obfuscation handling, procedure
      discipline, forensic reasoning, reporting quality.
  \end{itemize}
\end{frame}

% 9
\begin{frame}{Technical Implementation}
  \begin{itemize}
    \item \textbf{Frontend}: React and Vite.
    \item \textbf{Levels}: data-driven JSON dumps in \texttt{public/levels/}.
    \item \textbf{Generators}: deterministic Python scripts in \texttt{scripts/level\_generation/}.
    \item \textbf{Evidence}: source artefacts and UI visuals in \texttt{assets/};
      playable level data in \texttt{public/levels/}.
    \item \textbf{Media}: intro video in \texttt{public/intro.mp4}, with browser-safe audio enable.
    \item \textbf{Runtime}: browser-based interaction, no backend dependency during play.
  \end{itemize}

  \vspace{0.2cm}
  \begin{block}{Reproducibility}
    The campaign levels can be regenerated with \texttt{npm run levels:build}; the project
    can be built with \texttt{npm run build}. This keeps the game \textbf{data-driven}
    and reproducible.
  \end{block}
\end{frame}

% 10
\begin{frame}{Educational Value and Demo Plan}
  \begin{columns}[T,onlytextwidth]
    \begin{column}{0.49\textwidth}
      \begin{block}{Course alignment}
        \begin{itemize}
          \item \textbf{file signatures} and data carving;
          \item fragmentation and unallocated space;
          \item MBR and \textbf{Little-Endian} values;
          \item anti-forensics and decoys;
          \item partial recovery and defensible reporting.
        \end{itemize}
      \end{block}
    \end{column}
    \begin{column}{0.49\textwidth}
      \begin{block}{Suggested demo}
        Show the \textbf{evidence intake}, Act 1 for validation, Act 4 for offset
        reasoning, Act 5 or 6 for XOR, then close with the \textbf{timeline board}
        and final competency report.
      \end{block}
    \end{column}
  \end{columns}

  \vfill
  \centering
  \Large Questions?
\end{frame}

\end{document}
